\subsection{Algoritma String Matching}
\textbf{Definisi:} String matching adalah proses mencari kemunculan pola (\textit{pattern}) di dalam teks.

\textbf{Masukan:} Teks $T$ dengan panjang $n$ dan pola $P$ dengan panjang $m$.

\textbf{Keluaran:} Indeks awal kemunculan pola, atau -1 jika pola tidak ditemukan.

\subsection{Algoritma KMP}
\textbf{Definisi:} Knuth-Morris-Pratt (KMP) adalah algoritma pencocokan string dari kiri ke kanan.

\textbf{Cara Kerja:} KMP membuat tabel LPS untuk menentukan pergeseran pola saat terjadi ketidakcocokan.

\textbf{Kelebihan:} Tidak mengulang perbandingan karakter teks yang sudah diketahui.

\subsection{Algoritma Boyer-Moore}
\textbf{Definisi:} Boyer-Moore adalah algoritma pencocokan string dari kanan ke kiri.

\textbf{Cara Kerja:} Algoritma memakai tabel \textit{last occurrence} untuk menentukan besar pergeseran pola.

\textbf{Kelebihan:} Pada teks panjang, pergeseran pola dapat lebih jauh dibandingkan pencocokan biasa.
